% Technical appendix: methods, specifications, known limitations, and how to
% reproduce every number. Written so an outside reader can rerun the analysis.
\section{Technical appendix}
\label{sec:appendix}

This appendix records how to rerun the analysis, how the estimates and
regressions are specified, and where the evidence runs out. Every method in
this appendix names the Stata do-file that produced it, so a reader can open
the file, read the code, and reproduce the number. All do-files are in
\texttt{03\_analysis/stata/}, and the merged number ledger is at
\texttt{03\_analysis/out/canonical\_numbers.csv}.

\subsection{Reproducing this study}

The analysis folder is self-contained: a reader can copy it to another machine
and rerun it without editing a path, because every path resolves relative to
the project root, which the setup file \texttt{\_setup.do} locates by
searching upward for a known sentinel file. Outside data files are stored in
\texttt{03\_analysis/data/} rather than at referenced absolute paths.
Rebuilding the study from raw data takes three commands, run in order from
the analysis folder.

\begin{center}
{\footnotesize
\begin{tabular}{@{}B{4.2cm} P{12.6cm}@{}}
\toprule
\textbf{Step} & \textbf{Command} \\
\midrule
Run the descriptive modules & \texttt{cd 03\_analysis/stata \&\& stata-mp -b do run\_all.do} \\
\rowrule
Run the quasi-experimental module & \texttt{stata-mp -b do 70\_quasi\_experiments.do} \\
\rowrule
Build the study & \texttt{cd ../../05\_report \&\& make} \\
\bottomrule
\end{tabular}}
\end{center}

\texttt{run\_all.do} runs the seven descriptive modules that produce the
figures and registered numbers in Sections~\ref{sec:earnings} through
\ref{sec:state}. Those modules are, in order,
\texttt{10\_pums\_earnings.do} for PUMS earnings,
\texttt{20\_wages\_industry.do} for the payroll wage and industry series,
\texttt{30\_venues.do} for the venue panel,
\texttt{40\_city\_money.do} for the city funds,
\texttt{50\_state\_national.do} for the state and national comparison,
\texttt{60\_cost\_of\_living.do} for the cost-of-living series, and
\texttt{80\_census\_microdata.do} for the 2022 music census.
\texttt{run\_all.do} finishes by calling
\texttt{build\_numbers.do}, so the number macros need no separate command.
\texttt{run\_all.do} does not call the quasi-experimental module
\texttt{70\_quasi\_experiments.do}, which is the sole source of every quantity
behind the four designs in Sections~\ref{sec:venues} and \ref{sec:state}, so
that module runs separately as the second step. Because
\texttt{build\_numbers.do} merges whatever ledgers it finds on disk, skipping
the second step produces a report that still compiles, using stale
quasi-experimental numbers.

Every quantity registered in the numbers ledger is quoted through a generated
macro rather than a typed figure. A small number of figures quoted from named
source documents, and rounded verbal descriptions of registered quantities,
are typed in place and cited where they appear. Each analysis module writes
its results to its own ledger, and \texttt{build\_numbers.do} merges those
ledgers into the macro file the study reads. The merged ledger at
\texttt{03\_analysis/out/canonical\_numbers.csv} lists every registered
quantity with its value, its source file, and a \texttt{note} field recording
the population, filter, weight and vintage behind it. A reader checking any
number in this study can find it in that file and trace it to the module that
produced it.

\subsection{Inflation adjustment}

All dollar figures are real, in 2025 dollars, converted with the U.S. Bureau
of Labor Statistics Consumer Price Index for All Urban Consumers (series
CUUR0000SA0, national, all items), and every figure subtitle states the base
year.\footnote{Series documentation and monthly values at
\url{https://www.bls.gov/cpi/}; this study pulls the series from the BLS
Public Data API at \url{https://api.bls.gov/publicAPI/v2/timeseries/data/}.}
Three kinds of figure are flagged nominal where they appear: a ratio of two
same-year dollar figures, where deflation cancels; a figure quoted as its
source originally published it; and award or excluded-filing totals whose
dollars all fall inside one period.

The Bureau of Labor Statistics never published a Consumer Price Index value
for October 2025, because the federal funding lapse suspended
collection.\footnote{The Bureau paused the October 2025 CPI collection during
the fiscal-year 2026 government shutdown; the September and November values
are published as scheduled. See BLS release schedule at
\url{https://www.bls.gov/schedule/news_release/cpi.htm}.} This analysis
interpolates that month from September and November and flags it in the data,
so the 2025 annual average that deflates every dollar figure in this study
rests on twelve months rather than eleven. The interpolation runs inside
\texttt{60\_cost\_of\_living.do}.

Austin has no local price index, as Section~\ref{sec:measurement} notes.
Deflating by the national index is the only option available, and its effect
on the estimates here is unknown in sign. The one cross-sectional price
comparison that does exist puts Austin's all-items Regional Price Parity at
\NColRppAllitemsAustin\ against a national 100, which says nothing about the
rate at which Austin prices moved.\footnote{U.S. Bureau of Economic Analysis,
Regional Price Parities by Metropolitan Statistical Area, 2024 vintage, at
\url{https://www.bea.gov/data/prices-inflation/regional-price-parities-state-and-metro-area}.
The Regional Price Parity indexes the price level in a metro against a
national 100 in the same year, and its own documentation says it is not a
time-series measure of local inflation.}

\subsection{Survey estimates and their uncertainty}

Microdata estimates, computed from individual survey records rather than from
published tables, use the American Community Survey five-year Public Use
Microdata Sample (PUMS) for Texas, 2020 to 2024, the newest five-year PUMS
released as of retrieval.\footnote{U.S. Census Bureau, American Community
Survey 2020--2024 five-year PUMS for Texas, released 5 March 2026, at
\url{https://www2.census.gov/programs-surveys/acs/data/pums/2024/5-Year/csv_ptx.zip}.
The ACS is the Bureau's ongoing household survey, and the PUMS is its
individual-record release; it is the only public federal source that supports
occupation-by-metro estimates for a workforce as small as Musicians and
Singers because it pools the occupation across five years.} Dollar variables
are adjusted first by the survey's own \texttt{ADJINC} income-adjustment
factor, which places five survey years on a common basis, and then by the
Consumer Price Index.

Standard errors come from the Census Bureau's successive-difference
replication method, the published approach for this file, which re-estimates
each statistic on all 80 replicate weights (\texttt{PWGTP1} to
\texttt{PWGTP80}) and takes the spread across those estimates as the
error.\footnote{U.S. Census Bureau, \emph{PUMS Accuracy of the Data} technical
document for the 2020--2024 five-year file, at
\url{https://www2.census.gov/programs-surveys/acs/tech_docs/pums/accuracy/2020-2024AccuracyPUMS.pdf}.
Treating microdata records as an independent sample would understate the true
standard error substantially, because the ACS is a stratified cluster
sample rather than a simple random sample.} Margins of error are reported at
90 percent. The PUMS extractions and every survey estimate in this study run
in \texttt{10\_pums\_earnings.do}.

Two population definitions are used, and each is applied to both sides of
every comparison it appears in; Section~\ref{sec:formulas} defines them and
names which estimates use which. This study suppresses estimates resting on
fewer than 50 unweighted records, the practice the Bureau recommends. The
Austin musician cell has \NPumsAustinMusNPos\ unweighted person records,
which supports the headline estimates reported here and does not support any
further breakdown.

\subsection{Regression specifications}

\textbf{Conditional earnings gap.} This study regresses the natural logarithm
of personal earnings on an indicator for music occupations by weighted least
squares, controlling for age and age squared, sex, educational attainment,
usual hours worked, weeks worked and metro fixed effects, among employed
Texans aged 18 to 64 with positive earnings, with standard errors from the
80 replicate weights. The coefficient is \NPumsRegMusicCoef\ (standard error
\NPumsRegMusicSe), a difference of \NPumsRegMusicPctgap. The full table is
at \texttt{03\_analysis/out/tables/table\_earnings\_regression.csv} and the
regression runs in \texttt{10\_pums\_earnings.do}.

\textbf{Venue divergence.} A two-way fixed-effects regression of the log of
real monthly receipts on a core-live-music indicator interacted with the
post-March-2020 period, absorbing venue and year-month fixed effects and
clustering standard errors by venue, gives \NVenueDidCoef\ (standard error
\NVenueDidSe, p = \NVenueDidP) on \NVenueDidN\ venue-months. The event-study
version replaces the single interaction with year-specific terms, omitting
2019.\footnote{A full factorial interaction under two-way fixed effects
leaves the term of interest collinear and the estimator drops it, so this
study constructs the interaction as an explicit product variable. Requesting
2019 as the base year is likewise not honoured when the fixed effects absorb
the alternatives, so the year-by-treatment terms are built by hand with 2019
omitted. Both fixes change the coefficients, which is why they are recorded
here. Pre-2019 coefficients trend upward rather than staying flat, so the
parallel-paths assumption does not hold cleanly and the design is
descriptive.} The divergence regression runs in \texttt{30\_venues.do}.

\textbf{Quasi-experimental designs.} Four designs appear in
Sections~\ref{sec:venues} and \ref{sec:state}: an event study around federal
relief awards, a difference-in-differences on the 2017 extended-hours pilot,
an event study around the 2024 venue grants, and a distance-band test around
the 2022 arena opening. For the two designs with the smallest treated groups
(the extended-hours pilot and the 2024 venue grants), conventional
cluster-robust standard errors are unreliable, so this study uses
randomization inference as the primary test: the module reassigns treatment
at random across the eligible venues several thousand times, and the exact
p-value is the share of those placebo estimates at least as large as the
observed one. The arena and federal relief designs report Wald tests on
cluster-robust standard errors instead; the arena design's reference group
is six venues, so its bands are imprecisely estimated and none is
individually distinguishable from zero. All four designs are descriptive:
selection into treatment is not random in any of them, and the relief design
in particular selects on the severity of the shock by construction. All four
designs run in \texttt{70\_quasi\_experiments.do}. Section~\ref{sec:formulas}
below gathers the specification behind each design with every symbol defined.

\subsection{Every formula in one place}
\label{sec:formulas}

The footnotes at each point of use give the specification behind each
estimate. This subsection gathers all of them, so a reviewer can read the
methods in one sitting instead of hunting through the sections. Nothing here
is new: each block names the estimator, defines every symbol, names the
estimation sample and the variance assumption, and says whether the quantity
is descriptive or causal. Where a quantity appears in more than one section,
this appendix writes it once here.

\subsubsection*{Notation used throughout}

Subscript \(i\) indexes a unit: a person in the survey blocks and a venue or
permit location in the panel blocks. Subscript \(t\) indexes time: a month in
the venue panel and a year elsewhere. \(C_y\) is the annual average of the
national Consumer Price Index for All Urban Consumers in year \(y\), and
\(d_y = C_{2025}/C_y\) is the factor that carries year \(y\) dollars to this
study's 2025 base. \(w_i\) is a survey weight, \texttt{PWGTP} for persons and
\texttt{WGTP} for households in the American Community Survey PUMS.
\(\mathbf{1}(\cdot)\) equals one when the condition inside holds and zero
otherwise. A hat marks an estimate.

\subsubsection*{Price adjustment}

Nominal dollars become real dollars by \(x^{\mathrm{real}}_y = x_y \times
d_y\). The 2025 annual average uses an interpolated October, the average of
September and November, because the federal funding lapse suspended
collection that month. A ratio of two same-year dollar figures is left
nominal on purpose, because deflating both sides by the same \(d_y\) leaves
the ratio unchanged. The deflator ledger \texttt{cpi\_annual.dta} is built
inside \texttt{60\_cost\_of\_living.do} from the BLS Public Data API.

Microdata dollars take an extra step first. For person \(i\) the real value
is \(y_i = Y_i \times (\mathrm{ADJINC}_i/10^{6}) \times d_{2024}\), where
\(Y_i\) is the amount as reported and \(\mathrm{ADJINC}_i\) is the survey's
own income adjustment factor, published with six implied decimals, which
restates all five survey years on a common 2024
basis.\footnote{Documentation at
\url{https://www2.census.gov/programs-surveys/acs/tech_docs/pums/2020-2024_PUMS_README.pdf}
(ADJINC is section 3.3 of the readme).} Housing dollars use
\texttt{ADJHSG} in the same position. No year-specific deflator is merged
on, because the survey factor has already removed the inflation inside the
2020 to 2024 window and a second one would remove it twice.

\subsubsection*{Weighted descriptive estimates and their uncertainty}

A weighted median is
\(\hat\theta_{0.5} = \inf\{\,y : \sum_i w_i \mathbf{1}(y_i \le y)
   \ge 0.5 \sum_i w_i\,\}\),
the smallest observed value at which accumulated weight reaches half the
total, with adjacent order statistics averaged when the accumulated weight
lands exactly on half. A weighted share of a group \(G\) is
\(\hat p = \sum_{i \in G} w_i a_i \big/ \sum_{i \in G} w_i\), where \(a_i\)
is one when the record has the attribute being counted. Both are descriptive
population estimates.

Sampling error comes from successive-difference replication (SDR), the
published method for this file:
\[
\widehat{\operatorname{Var}}(\hat\theta)
  = \frac{4}{80}\sum_{r=1}^{80}\bigl(\hat\theta_r - \hat\theta\bigr)^{2},
\qquad
\mathrm{MOE}_{90} = 1.645\,\sqrt{\widehat{\operatorname{Var}}(\hat\theta)},
\]
where \(\hat\theta\) uses the full weight and \(\hat\theta_r\) the \(r\)th
replicate weight, \texttt{PWGTP1} to \texttt{PWGTP80} for persons and
\texttt{WGTP1} to \texttt{WGTP80} for households.\footnote{The constant
\(4/80\) is fixed by the Bureau's replicate design; the derivation is in
\emph{PUMS Accuracy of the Data}, section 5, cited above.} Deviations are
taken around the full-sample estimate rather than the replicate mean, the
convention Stata labels \texttt{mse}. A few replicate weight cells are
slightly negative, which the Bureau permits; this study floors the weight
at zero inside percentile calculations only, because a weighted percentile
is undefined with negative weights. For a difference between two
independent estimates the combined margin is
\(\mathrm{MOE}_{A-B} = \sqrt{\mathrm{MOE}_A^{2} + \mathrm{MOE}_B^{2}}\).

Two population definitions are used, and each is applied to both sides of
every comparison it appears in. The employed base is Texas residents aged
18 to 64 in work (employment status \texttt{ESR} in \{1,2,4,5\}). The
positive-earnings base adds the requirement that personal earnings exceed
zero. Medians, the self-employment share, coverage rates and the regression
run on the bases named in their own footnotes; the earnings-threshold
shares run on the employed base, so that people reporting zero or negative
earnings stay in the denominator. No estimate resting on fewer than 50
unweighted records is published. All weighted descriptive estimates run in
\texttt{10\_pums\_earnings.do}.

\subsubsection*{Conditional earnings gap}

\[
\ln y_i = \alpha + \beta\,\mathrm{music}_i
  + \gamma_1 \mathrm{age}_i + \gamma_2 \mathrm{age}_i^{2}
  + \gamma_3 \mathrm{female}_i
  + \sum_{k=2}^{5}\delta_k\,\mathrm{educ}_{ki}
  + \gamma_4 \mathrm{hours}_i + \gamma_5 \mathrm{weeks}_i
  + \sum_{m}\lambda_m\,\mathrm{metro}_{mi} + \varepsilon_i
\]

Every symbol appears in the footnote at the point of use in
Section~\ref{sec:earnings}. Estimation is by weighted least squares with the
person weight, over the \NPumsRegN\ employed Texans aged 18 to 64 with
positive earnings, with standard errors from the 80 replicate weights
rather than from an independence assumption. A log-point coefficient
converts to a percentage difference as \(100\,(e^{\hat\beta} - 1)\), and the
same transformation applied to \(\hat\beta \pm 1.96\,\mathrm{se}(\hat\beta)\)
gives the 95 percent interval. Two nested specifications, metro fixed
effects alone and metro plus demographics, appear in the exported table
beside the one quoted in the text. The estimate is a descriptive conditional
difference, not the earnings cost of choosing music: hours and weeks are
themselves outcomes of the occupation, so controlling for them removes part
of what is being measured and pushes the coefficient toward zero. The
regression runs in \texttt{10\_pums\_earnings.do}.

\subsubsection*{Index numbers, ratios and per-resident measures}

An index rebases a series on a chosen year or month,
\(I_{s,t} = 100 \times X_{s,t}/X_{s,t_0}\). Three bases appear in this
study, each chosen for what its figure compares: each series against its
own first available year, for payroll jobs against nonemployer businesses;
each metro against its own January 2015, for real rent; and every group
against 2019, for venue receipts. An index carries no unit and no count,
which is why two indexed panels can be compared for direction and never
added together. The index constructions run in
\texttt{20\_wages\_industry.do} (jobs against nonemployer),
\texttt{60\_cost\_of\_living.do} (metro rent) and \texttt{30\_venues.do}
(receipts).

An endpoint percentage change is \(100\,(X_T/X_0 - 1)\), a total change
across the whole span rather than an annual rate.

The hours-of-work ratio is \(h_{o,y} = F_y / w_{o,y}\), monthly Fair Market
Rent over the median hourly wage for occupation group \(o\), matched within
year and left in nominal dollars, because
\((d_y F_y)/(d_y w_{o,y}) = F_y / w_{o,y}\). Fair Market Rent is the U.S.
Department of Housing and Urban Development benchmark rent used to set
voucher payment standards, published annually by HUD metro area at
\url{https://www.huduser.gov/portal/datasets/fmr.html}. The wages come from
the BLS Occupational Employment and Wage Statistics program, described
below. The ratio is computed in \texttt{60\_cost\_of\_living.do}.

The rent share of earnings is \(s = 100 \times 12\bar{F}/\tilde{y}\), the
annualised average Fair Market Rent over the five fiscal years the microdata
pool, divided by the weighted median earnings of Austin musicians, both in
2025 dollars.

Businesses per 10,000 residents is \(10^{4} \times E_{c,y}/N_{c,y}\), and
receipts per business is \(\bar{r}_{c,y} = R^{\mathrm{real}}_{c,y}/E_{c,y}\),
an arithmetic mean with no distribution published behind it. Both use the
Census Bureau's Nonemployer Statistics for the numerator, which reports
counts of establishments with no paid employees and no covered payroll, and
the Bureau's Vintage 2025 population estimates for the
denominator.\footnote{Nonemployer Statistics documentation at
\url{https://www.census.gov/programs-surveys/nonemployer-statistics.html};
Vintage 2025 state and county estimates at
\url{https://www.census.gov/programs-surveys/popest.html}. This branch of
the analysis runs in \texttt{20\_wages\_industry.do}.}

Per-resident state support is computed twice from the same denominator,
\(a_s = A_s\,d_{y(A_s)}/N_s\) for the authorised amount and
\(b_s = B_s\,d_{y(B_s)}/N_s\) for the amount disbursed, where \(N_s\) is
resident population on 1 July 2025 from the Census Bureau's Vintage 2025
national estimates. An authorisation is a statutory ceiling and a
disbursement is money that left the treasury, and the two can differ by the
whole amount. Both computations run in \texttt{50\_state\_national.do}.

Regional Price Parity gaps are
\(g_m = \mathrm{RPP}^{\mathrm{hous}}_m - \mathrm{RPP}^{\mathrm{all}}_m\),
and Consumer Price Index gaps are differences between two series rebased on
the same year. Both are in index points and neither is a percentage. Both
are computed in \texttt{60\_cost\_of\_living.do}.

Every measure in this block is an arithmetic ratio of published quantities,
so none carries a sampling error of its own, though the rent share inherits
the margin of error of the earnings median in its denominator.

\subsubsection*{Panel regressions on venue receipts}

The curated panel of \NVenuePanelIsCurated\ Austin live-music rooms supplies
the treated group. The comparison group in the divergence and event-study
designs is every other mixed-beverage permit location in Travis County and
the City of Austin, which is why those two regressions run on \NVenueDidN\
venue-months rather than on the panel's own \NQuasiPanelVenueMonths. The
four quasi-experimental designs run inside the panel and its neighbours
instead.

The venue designs below share an outcome and a fixed-effects structure. The
outcome is \(\ln R^{\mathrm{real}}_{it}\), the log of venue \(i\)'s
mixed-beverage receipts in month \(t\) in 2025 dollars, taken from Texas
Comptroller Mixed Beverage Gross Receipts (dataset \texttt{naix-2893}) at
\url{https://data.texas.gov/dataset/Mixed-Beverage-Gross-Receipts/naix-2893}.
\(\alpha_i\) is a venue fixed effect absorbing everything constant about the
room, and \(\tau_t\) is a year-month fixed effect absorbing everything
common to the market in that month. Standard errors are clustered by venue
throughout, which allows one room's months to be correlated with one another
and treats venues as the independent units. Zero-receipt months leave a log
specification, so every log estimate describes venues conditional on
trading that month. All venue regressions run in \texttt{30\_venues.do}
(divergence and event study) and \texttt{70\_quasi\_experiments.do}
(federal relief, extended-hours pilot, 2024 venue grants and arena
opening).

\begin{itemize}
\item \textbf{Divergence, core tier against the rest of the market.}
      \(\ln R^{\mathrm{real}}_{it} = \alpha_i + \tau_t
        + \beta(\mathrm{Core}_i \times \mathrm{Post}_t) + \varepsilon_{it}\),
      with \(\mathrm{Post}_t\) equal to one from March 2020, over
      \NVenueDidN\ from January 2013 to December 2025. The comparison group
      is other mixed-beverage permit locations in Travis County and Austin;
      the panel's live-music-adjacent tier is dropped from the estimation
      sample rather than coded zero, because those rooms programme live
      music intermittently and belong to neither group.
\item \textbf{Event study.} The same equation with
      \(\sum_{y \neq 2019}\beta_y(\mathbf{1}\{\mathrm{year}(t)=y\}
        \times \mathrm{Core}_i)\) in place of the single interaction.
\item \textbf{Federal relief.} Half-year blocks relative to each venue's own
      award month, \(\mathrm{block}_{it} = \lfloor (t-a_i)/6 \rfloor\),
      clipped at nine blocks before and eight after, with the block 30 to
      25 months before the award omitted. A second version replaces the
      recipient indicator with award intensity, the award divided by mean
      monthly 2019 receipts at the same venue, so one unit is one month of
      pre-pandemic revenue. Award data come from the U.S. Small Business
      Administration's Shuttered Venue Operators Grant awards file (awards
      as of 5 July 2022) at
      \url{https://data.sba.gov/sites/default/files/distribution/SBA-ODA-2022-09-001/awards-as-of-7-5-22.xlsx}.
\item \textbf{Extended-hours pilot.}
      \(\mathrm{Pilot}_i \times W_t\) over January 2015 to December 2019,
      with a half-year event study on a 2016 second-half base and a variant
      adding venue-by-month-of-year fixed effects, which absorb the
      seasonality the two groups do not share. \(W_t\) equals one during
      the pilot's operating window on Red River.
\item \textbf{Venue grants.} \(\mathrm{Grant}_i \times \mathrm{Post2024}_t\)
      over January 2018 to December 2025, with an annual event study on a
      2023 base. The grant set comes from the City of Austin's fiscal 2024
      Live Music Fund awardee list at
      \url{https://austin.widen.net/content/bpzuacfkk8/pdf/FY24-AwardeeList_Austin-Live-Music-Fund.pdf}.
\item \textbf{Arena opening.} Three distance-band interactions with the
      post-April-2022 period, venues more than four miles away omitted, and
      the pandemic window dropped from both sides. Straight-line distance
      comes from the great-circle formula
      \[
      D = 2\rho\arcsin\sqrt{\sin^{2}(\Delta\phi/2)
        + \cos\phi_1\cos\phi_2\sin^{2}(\Delta\lambda/2)},
      \]
      with \(\phi\) latitude, \(\lambda\) longitude, \(\Delta\) the
      difference between venue and arena, and \(\rho\) the Earth's radius
      in miles.
\end{itemize}

Interactions are constructed as explicit product variables and base periods
are dropped by hand throughout. Under two-way fixed effects a factorial
interaction leaves the term of interest collinear and the estimator drops
it, and a requested base year is not honoured once the fixed effects
absorb the alternatives, which renormalises an event study onto whichever
level the estimator happened to drop.

The exported venue table carries a fourth column that replaces the log with
the inverse hyperbolic sine,
\[
\operatorname{asinh}(R) = \ln\bigl(R + \sqrt{R^{2}+1}\bigr),
\]
which behaves like a logarithm at large values and is defined at zero, so
months when a room reported no receipts stay in the sample. That column
answers a different question from the log column, which conditions on
trading, and its coefficient is not scale-free: rescaling the currency
changes its size, so a reader cannot read it as a percentage. No number
quoted in this study comes from it, and the column stays in the table
because dropping zero months is the log specification's main limitation and
a reader should be able to see what happens without that drop.

Joint tests reported as p-values are Wald tests. The parallel-paths checks
test that the pre-treatment event-study coefficients are jointly zero, and
the arena test is \(H_0: \beta_1 = \beta_2 = \beta_3\), which asks whether
the pattern varies with distance at all rather than whether any single band
moved.

Every venue design is descriptive. Nothing in them is randomised, no policy
switches on for one group and not another at a clean date, and the relief
design selects on the severity of the shock by construction.

\subsubsection*{Randomization inference}

Where the treated group is small, cluster-robust standard errors understate
the true sampling variability, so the exact p-value comes from permuting
the treatment assignment:
\[
p = \frac{1 + \#\{\,r \le R : |\hat\beta_r| \ge |\hat\beta|\,\}}{R + 1},
\]
where \(\hat\beta\) is the estimate under the assignment that actually
happened, \(\hat\beta_r\) is the estimate from refitting the same
regression after reassigning treatment at random to the same number of
units, drawn without replacement from an eligible pool, and \(R\) is the
number of draws: 5{,}000 for the extended-hours pilot and 2{,}000 for the
venue grants. Treatment dates are held fixed and only the assignment moves.
Adding one to numerator and denominator counts the assignment that happened
among the possibilities, which is what keeps the test exact rather than
approximate.\footnote{The derivation is in Fisher's \emph{Design of
Experiments}, 1935, chapter 2, and the small-sample properties of the
adjustment are in Phipson and Smyth, ``Permutation p-values should never be
zero,'' \emph{Statistical Applications in Genetics and Molecular Biology}
9(1), 2010.} The pool is restricted to venues whose filing coverage matches
the treated group's, so a placebo assignment lands on the same kind of unit
as the real one. The spread of the placebo estimates also bounds what the
design could have detected, and the report quotes that bound where it is
informative. The randomization loop runs in
\texttt{70\_quasi\_experiments.do}.

\subsubsection*{Association tests on the 2022 census file}

The 2022 Greater Austin Music Census carries no design weights, because it
is a self-selected online instrument fielded through community mailing
lists between 15 July and 12 September 2022 and not a probability sample,
so every statistic in that section is an unweighted count over a stated
denominator, \(\hat p = 100 \times n_{\mathrm{yes}}/n_{\mathrm{answered}}\),
and no margin of error exists for any of them.\footnote{Sound Music Cities
and the City of Austin Music \& Entertainment Division, \emph{2022 Greater
Austin Music Census}, summary report at
\url{https://www.austinmusiccensus.org/s/2022-Greater-Austin-Music-Census-SUMMARY-REPORT-v13-optimized.pdf};
the respondent-level microdata are published on the City of Austin's open
data portal as dataset \texttt{an3p-3yqx} at \url{https://data.austintexas.gov/}.} The
census-file estimates all run in \texttt{80\_census\_microdata.do}.

Independence between two survey items is tested with Pearson's chi-squared,
\(X^{2} = \sum_j \sum_k (O_{jk}-E_{jk})^{2}/E_{jk}\) with
\(E_{jk} = R_j C_k/N\), on \((J-1)(K-1)\) degrees of freedom, where
\(O_{jk}\) is the observed count in row \(j\) and column \(k\), \(R_j\) and
\(C_k\) are the row and column totals, and \(N\) is the number answering
both items.

The retention regression is a logit,
\[
\ln\frac{\Pr(\mathrm{stay}_i = 1)}{1-\Pr(\mathrm{stay}_i = 1)}
  = \alpha + \sum_s \beta_s\,\mathrm{sector}_{si}
  + \sum_e \gamma_e\,\mathrm{exper}_{ei}
  + \sum_h \delta_h\,\mathrm{tenure}_{hi}
  + \theta\,\mathrm{insured}_i ,
\]
estimated by maximum likelihood on the \NCensusLogitN\ respondents who
answered every item in the equation, with the quoted quantity the odds
ratio \(e^{\hat\delta_{\mathrm{renter}}}\). An odds ratio is not a ratio of
probabilities and is further from one than the probability ratio is
when the outcome is common, as it is here. Three nested models appear in
the exported table. The logit runs in \texttt{80\_census\_microdata.do}.

The published spending headline and the per-respondent total are two
different calculations. The headline is \(\sum_{k=1}^{8}\bar{x}_k\), a sum
of eight category means each taken over the respondents who answered that
category. The per-respondent figure is
\(\tilde{x} = \operatorname{median}_{i \in S}\bigl(\sum_{k=1}^{8}x_{ik}\bigr)\)
over \(S = \bigcap_k S_k\), the respondents who answered every category,
so each figure in the sum belongs to one person.

Both the tests and the logit assume a simple random sample, which a
self-selected online survey is not, so they describe which characteristics
separate respondents from one another and support no inference to the
Austin music workforce.

\subsection{The venue panel}

The panel matches Texas Comptroller Mixed Beverage Gross Receipts (dataset
\texttt{naix-2893}) to named Austin venues by location address and taxpayer
name, following venues through changes of owner and legal name, and covers
\NVenuePanelIsCurated\ rooms monthly from January 2007 to June
2026.\footnote{Comptroller portal at
\url{https://data.texas.gov/dataset/Mixed-Beverage-Gross-Receipts/naix-2893};
the file records every mixed-beverage permit holder's monthly gross
receipts, cover-charge receipts (unreliable before roughly 2022 because of
a reporting change), and permit status. The matching runs in
\texttt{30\_venues.do} and uses a curated venue list at
\texttt{01\_evidence/08\_venues\_ecosystem/venue\_match\_list.csv}.} Four
limitations bound what the panel can show.

\begin{itemize}
\item It measures alcohol sales, not music activity, and nothing in it
      records what performers were paid.
\item It is curated, not a census, and over-represents nameable rooms in
      the central city, particularly Red River, East Austin and South
      Congress.
\item It omits venues without a liquor permit, which removes several
      well-known listening rooms (Cactus Cafe, Meanwhile Brewing, Carousel
      Lounge, Bass Concert Hall, Central Presbyterian) from the panel.
\item Festival companies file under venue-linked names. The largest such
      filing is more than three hundred times its venue's median month, and
      the ten excluded venue-months are worth
      \NVenueFestivalDollarsExcluded\ in nominal dollars in total.\footnote{Two modules
      apply slightly different exclusion rules, one dropping the flagged
      extreme filings and the other a wider set. Testing the difference
      showed it does not matter: the core-tier change from 2019 to 2025 is
      identical under both rules. The cover-charge field is separately
      unusable before about 2022 because of a reporting change, so this
      study does not use it.}
\end{itemize}

\subsection{Known limitations}

\begin{itemize}
\item \textbf{No current representative measurement of Austin musician
      earnings exists.} The survey estimates in this study are the best
      available and rest on \NPumsAustinMusNPos\ Austin records pooled over
      five years of the American Community Survey PUMS.
\item \textbf{The federal wage survey excludes the self-employed.} The
      Bureau of Labor Statistics Occupational Employment and Wage
      Statistics (OEWS) program is a semi-annual establishment survey; it
      records the wage of every worker on a respondent establishment's
      payroll and records nothing for the self-employed, who fall outside
      the sampling frame.\footnote{OEWS program documentation at
      \url{https://www.bls.gov/oes/oes_ques.htm}. The Bureau also does not
      publish an annual musician wage for any city or year because the work
      is too irregular to annualise from an hourly rate.} Self-employed
      workers are \NPumsAustinMusSelfemp\ (\NPumsAustinMusSelfempMoe) of
      Austin musicians in the ACS PUMS, so even at the low end of that
      margin every figure drawn from OEWS describes a minority of the
      occupation.
\item \textbf{The 2022 Greater Austin Music Census has no design weights}
      and is a self-selected online sample, so its percentages describe
      respondents only and have no margins of error; a probability-sample
      benchmark for the Austin music workforce does not exist.
\item \textbf{The Arts and Cultural Production Satellite Account has no
      metro detail and will not be updated after 2023.} The Bureau of
      Economic Analysis announced in February 2026 that the April 2025
      release is the final regular annual release.\footnote{BEA news
      release 25-13, at
      \url{https://www.bea.gov/data/special-topics/arts-and-culture}.}
\item \textbf{Metro groupings for federal Shuttered Venue Operators Grant
      relief} are built from recipient city names rather than official
      metropolitan-statistical-area boundaries and are approximate.
\item \textbf{One award document does not reconcile with itself in three
      places}, and a second disagrees with press reporting by one award.
      This study uses each document's stated headline totals rather than
      re-tallying its line items.
\item \textbf{Records the city and the state hold but do not publish}:
      applicant demographics for the Austin Live Music Fund, the recipient
      list for the Texas Music Incubator Rebate Program, and the fund's
      audited fund-balance schedule. Section~\ref{sec:next} sets out how a
      requester could obtain them under the Texas Public Information Act
      (Texas Government Code Chapter 552) or the City of Austin's
      equivalent process.
\end{itemize}

\subsection{Software}

The analysis runs in Stata 19 and uses these user-written commands:
\texttt{reghdfe} and \texttt{ftools} for high-dimensional fixed effects,
\texttt{estout} for tables, \texttt{spmap} and \texttt{shp2dta} for the map,
and \texttt{gtools}, \texttt{palettes}, \texttt{colrspace} and
\texttt{grstyle} for data handling and figure styling. The project bundles
all of them in \texttt{03\_analysis/stata/ado/}, so no package installation
is needed, and \XeLaTeX{} typesets the document.
